\DocumentMetadata{
  tagging=on,
  tagging-setup={math/setup=mathml-SE,math/alt/use} 
}
\documentclass[12pt,letterpaper]{article}
\usepackage{fullpage}
\usepackage[top=2cm, bottom=4.5cm, left=2.5cm, right=2.5cm]{geometry}
\usepackage{amsmath,amsthm,amsfonts,amssymb,amscd,bm}
\usepackage{lastpage}
\usepackage{enumerate,multicol}
\usepackage{fancyhdr}
\usepackage{mathrsfs}
\usepackage{xcolor}
\usepackage{graphicx}
\usepackage{listings}
\usepackage{hyperref}
\usepackage{tabularx}


\usepackage{pgfplots}
\pgfplotsset{width=10cm,compat=1.9}
% \usepgfplotslibrary{external}
% \tikzexternalize
\usetikzlibrary{calc,angles,quotes,arrows,decorations.markings}

\definecolor{UTorange}{RGB}{207, 83, 0}
\definecolor{UCLAblue}{RGB}{39, 116, 174} 
\definecolor{UCLAdark}{RGB}{0, 59, 92}
\definecolor{UCLAlight}{RGB}{218, 235, 254}
\definecolor{UCLAgold}{RGB}{255, 209, 0}
\hypersetup{colorlinks,breaklinks,linkcolor=UCLAblue,urlcolor=UCLAdark,anchorcolor=UCLAblue,citecolor=black}
 
\newcommand{\R}{\mathbb{R}}
\newcommand{\C}{\mathbb{C}}
\newcommand{\Z}{\mathbb{Z}}
\newcommand{\Q}{\mathbb{Q}}
\newcommand{\F}{\mathbb{F}}
\newcommand{\N}{\mathbb{N}}

\newcommand{\inv}{^{-1}}

\newcommand{\ds}{\displaystyle}
\newcommand{\dydx}{\frac{dy}{dx}}
\newcommand{\dxdy}{\frac{dx}{dy}}
\newcommand{\dsint}{\displaystyle\int}


\newtheorem{theorem}{Theorem}[section]
\newtheorem{lemma}[theorem]{Lemma}

\theoremstyle{definition}
\newtheorem{definition}[theorem]{Definition}
\newtheorem{xca}[theorem]{Exercise}
\newtheorem{proposition}[theorem]{Proposition}
\newtheorem{corollary}[theorem]{Corollary}

\theoremstyle{remark}
\newtheorem{remark}[theorem]{Remark}


\setlength{\parindent}{0.0in}
\setlength{\parskip}{0.05in}

% Edit these as appropriate
\newcommand\course{Math 102}
\newcommand\hwnumber{1}                  % <-- homework number
\newcommand\name{netid19823}           % <-- NetID of person #1
\newcommand\NetIDb{netid12038}           % <-- NetID of person #2 (Comment this line out for problem sets)

\pagestyle{fancyplain}
\headheight 35pt
\lhead{\textbf{\Large \course}}
\chead{\textbf{\Large Review}}
\rhead{\textbf{\Large Jan 14, 2025}}
\lfoot{}
\cfoot{}
\rfoot{}
\headsep 1.5em

\begin{document}


%%%%%%%%%%%%%%%%%%%%
%accessibility/tagging
\hypersetup{
pdftitle={Math 102 Mad Minute Integration},
pdfauthor={Richard Wong},pdfdisplaydoctitle%
}
%%%%%%%%%%%%%%%%%%%%%%%

\begin{center}
\textbf{Name:} \underline{\hspace{2.2in}} \hfill \textbf{Net ID:} \underline{\hspace{2in}}

\end{center}


\begin{center}
    \fbox{You can submit this worksheet to Gradescope within 24 hours for 1 participation point.}
\end{center}
\subsection*{``Mad Minute" integration}

\noindent Antidifferentiate as many of the following functions as you can in two minutes. You may, for the purposes of this assignment, omit the added constant, $C$. 

\begin{center}

\tagpdfsetup{table/header-rows={1}}
\begin{tabularx}{\textwidth}{|c|X|}
  \hline \textbf{Function} & \textbf{Antiderivative} \\
\hline
\hline & \\ 
 $\sin(x)$  & \\ \hline
   \hline & \\ $   e^x$ & \\ \hline 
   \hline &\\$   x^3$ & \\ \hline 
   \hline &\\$   x^{-1}$ & \\ \hline 
   \hline &\\ $  (1+x^2)^{-1} $ & \\ \hline \hline &\\ $ \sec^2(x)$ & \\ \hline \hline &\\
  $ 2^x$ & \\ \hline \hline &\\
  $ e$ & \\ \hline \hline & \\
  $ \sec(x)\tan(x)$ & \\ \hline \hline &\\
  $ (1-x^2)^{-1/2}$ & \\ \hline \hline &\\
  $ \csc^2(x)$ & \\ \hline \hline &\\
  $ x(2+x)$ & \\ \hline \hline &\\
  $ (x-2)^{-1}$ & \\ \hline \hline & \\
  $ \csc(x)\cot(x)$ & \\ \hline 
\end{tabularx}
\end{center}




\end{document}
